\subsection{Towards de Sitter from 10D --- arXiv:1707.08678}

\textbf{Authors:} Jakob Moritz, Ander Retolaza, Alexander Westphal

\paragraph{主な主張}
非摂動的体積安定化の10次元記述に基づき、考察した単一モジュラスのKKLT構成ではupliftのbackreactionにより真空がAdSに留まる一方、racetrack安定化では回避できる可能性を論じる\cite{Moritz:2017xto}。

\paragraph{新規性}
非摂動効果とupliftの相互作用を10次元から検討し、四次元記述で想定するbackreactionの大きさを問題にする。

\paragraph{位置付け}
de Sitter upliftの10次元的整合性をめぐる論争の出発点の一つであり、後続の反論と併読すべき研究。

\paragraph{コメント（読書メモ）}
文献\cite{Kallosh:2019axr}を通じた読書メモでは、no-de-Sitter予想を支持する議論として位置付けた。ある版のKKLT構成では、10次元解析とbackreactionを含む四次元解析から、超対称AdS真空を準安定de Sitter真空へ整合的にupliftできないと主張している。一方、同文献の批判では、複雑な10次元解析に置かれた仮定や凝縮項の扱いに誤り・不足があり、再検討が必要だとされる。
